
\documentclass{article}
\usepackage{colortbl}
\usepackage{makecell}
\usepackage{multirow}
\usepackage{supertabular}

\begin{document}

\newcounter{utterance}

\centering \large Interaction Transcript for game `wordle\_withclue', experiment `medium\_frequency\_words\_clue\_no\_critic', episode 10 with reflexion{-}Qwen3.5{-}4B.
\vspace{24pt}

{ \footnotesize  \setcounter{utterance}{1}
\setlength{\tabcolsep}{0pt}
\begin{supertabular}{c@{$\;$}|p{.15\linewidth}@{}p{.15\linewidth}p{.15\linewidth}p{.15\linewidth}p{.15\linewidth}p{.15\linewidth}}
    \# & \multicolumn{2}{c}{Player} && \multicolumn{2}{c}{Game Master} \\
    \hline

    \theutterance \stepcounter{utterance}  
    & & & \multicolumn{4}{p{0.6\linewidth}}{
        \cellcolor[rgb]{0.9,0.9,0.9}{
            \makecell[{{p{\linewidth}}}]{
                \texttt{\tiny{[P1$\langle$GM]}}
                \texttt{You have attempted this same task before and did not succeed. These are the plans you wrote after those attempts.} \\
\texttt{Trial \#1: New plan: The last attempt was aborted because the explanation failed to explicitly state the final guess word, even though the word "tithing" was identified in the text. The rules require the format "explanation: details" followed immediately by "guess: word". I will ensure that after analyzing the feedback and clue to deduce the missing letter, I explicitly write "guess: {[}deduced word{]}" in the final line of my response, ensuring the word is clearly separated by the "guess:" tag.} \\
\\ 
\texttt{Trial \#2: New plan: The last attempt was aborted because the response format was incorrect; the "guess:" tag was missing or improperly placed, and the explanation did not strictly follow the required "explanation: details" followed immediately by "guess: word" structure. I will ensure that my response contains exactly two lines: the first line starting with "explanation:" containing the reasoning that combines the clue and any feedback, and the second line starting with "guess:" containing the five{-}letter target word. I will verify that the word is a valid five{-}letter English word and that the format matches the example precisely before sending.} \\
\\ 
\texttt{Trial \#3: New plan: The last attempt was aborted because the response format violated the strict requirement of having exactly two lines: the first starting with "explanation:" and the second starting with "guess:". Although I identified the word "tithing" from the clue, I likely failed to structure the output precisely as "explanation: {[}reasoning{]}" followed immediately by "guess: {[}word{]}". Next time, I will strictly adhere to the two{-}line format, ensuring the explanation clearly links the clue "tithing" to the target word, and the second line is exactly "guess: tithing", with no extra text, newlines, or deviations from the specified structure.} \\
\\ 
\\ 
\texttt{You are a language wizard who likes to guess words by using the given rules.} \\
\\ 
\texttt{Welcome to Wordle! You have six attempts to guess the target word, a valid English word of five lowercase letters (a{-}z). Please use the tags "explanation:" and "guess:" to provide a concise explanation for each guess.} \\
\\ 
\texttt{To help you make an informed guess, you will receive a clue for the word, such as} \\
\texttt{clue: snowy white.} \\
\\ 
\texttt{Here is an example guess based on the clue:} \\
\texttt{explanation: In the fairy tail Snow White, the girl is killed because she eats a poisoned apple. And the word apple has 5 letters.} \\
\texttt{guess: apple} \\
\\ 
\texttt{After each guess, your answer will be validated, and you will receive feedback indicating which letters are correct (green), which letters are correct but in the wrong position (yellow), and which letters are incorrect (red). This feedback can be useful in determining which letters to include or exclude in your next guess.} \\
\\ 
\texttt{For example, the feedback for "apple" might be:} \\
\texttt{guess\_feedback: a<yellow> p<yellow> p<green> l<yellow> e<red>} \\
\\ 
\texttt{The explanation should contain details about how both the guess\_feedback and clue work together in tandem to formulate a new guess.} \\
\\ 
\texttt{Do not generate any other text.} \\
\\ 
\texttt{Let's begin with your first guess.} \\
\\ 
\texttt{clue: tithing} \\
            }
        }
    }
     \\ \\

    \theutterance \stepcounter{utterance}  
    & \multicolumn{4}{p{0.6\linewidth}}{
        \cellcolor[rgb]{0.561,0.737,0.561}{
            \makecell[{{p{\linewidth}}}]{
                \texttt{\tiny{[P1$\rangle$GM]}}
                \texttt{explanation: The clue "tithing" refers to a voluntary contribution or offering, often religious, which is a synonym for "tithe". The word "tithe" fits the five{-}letter requirement and matches the context of the clue.} \\
\texttt{guess: tithe} \\
            }
        }
    }
    & & \\ \\

    \theutterance \stepcounter{utterance}  
    & & & \multicolumn{2}{p{0.3\linewidth}}{
        \cellcolor[rgb]{0.9,0.9,0.9}{
            \makecell[{{p{\linewidth}}}]{
                \texttt{\tiny{[GM$|$GM]}}
                \texttt{attempts = 1} \\
\texttt{target\_word = tenth} \\
\texttt{target\_word\_clue = tithing} \\
\texttt{guess = tithe} \\
\texttt{guess\_feedback = t<green> i<red> t<yellow> h<yellow> e<yellow>} \\
            }
        }
    }
    & & \\ \\

    \theutterance \stepcounter{utterance}  
    & & & \multicolumn{4}{p{0.6\linewidth}}{
        \cellcolor[rgb]{0.9,0.9,0.9}{
            \makecell[{{p{\linewidth}}}]{
                \texttt{\tiny{[P1$\langle$GM]}}
                \texttt{guess\_feedback: t<green> i<red> t<yellow> h<yellow> e<yellow>} \\
\\ 
\texttt{Provide your response only in this format.} \\
\texttt{explanation: details} \\
\texttt{guess: word} \\
            }
        }
    }
     \\ \\

    \theutterance \stepcounter{utterance}  
    & \multicolumn{4}{p{0.6\linewidth}}{
        \cellcolor[rgb]{0.561,0.737,0.561}{
            \makecell[{{p{\linewidth}}}]{
                \texttt{\tiny{[P1$\rangle$GM]}}
                \texttt{explanation: The feedback shows 't' is correct and in the first position (green), while 'i' is incorrect (red). The second 't' is in the wrong position (yellow), meaning it exists in the word but not in the third spot. The letters 'h' and 'e' are also in the wrong positions. This confirms the word starts with 't' and contains 'h' and 'e', but the pattern is not "t{-}i{-}t{-}h{-}e". Based on the clue "tithing" and the letter constraints, the word is likely "theft" or "thief", but given the clue context of religious giving, "tithe" was the initial guess. However, the feedback contradicts "tithe" because the second letter 'i' is red, not green. Re{-}evaluating the clue "tithing" as a hint for the word itself, perhaps the target is a synonym or related term. Wait, the clue is "tithing", which is a verb/noun. The word "tithing" has 7 letters. Maybe the clue is describing the word. Let's reconsider common 5{-}letter words related to tithing. "Tithe" was my first guess. The feedback says t<green> i<red> t<yellow> h<yellow> e<yellow>. This means the word has t, h, e, but not i. And the t is not in the 3rd spot. Could the word be "thief"? No, clue doesn't fit. "theft"? No. What if the clue is actually the word "tithing" and I need to find a 5{-}letter word that is a synonym? "Tithe" is the only direct synonym. Is it possible the feedback implies the word is "throe"? No. Let's look at the letters again: t, h, e are present. i is absent. The structure is t \_ \_ \_ \_. With h and e in wrong spots. Maybe "thine"? No. "thole"? No. Perhaps the clue is misleading or I am misinterpreting. Let's try "thief" again? No. What about "thorn"? No. Let's go back to the most likely candidate: "tithe". If the feedback is t<green> i<red> t<yellow> h<yellow> e<yellow>, then the word is NOT "tithe". But "tithe" is the only 5{-}letter word for tithing. Is there another word? "Tithes" is 6. "Tithing" is 7. Maybe the word is "tithy"? No. Let's assume the word is "thief" and the clue is just a distractor? No, that's unlikely. Let's try "throe"? No. How about "thole"? No} \\
            }
        }
    }
    & & \\ \\

    \theutterance \stepcounter{utterance}  
    & & & \multicolumn{2}{p{0.3\linewidth}}{
        \cellcolor[rgb]{0.9,0.9,0.9}{
            \makecell[{{p{\linewidth}}}]{
                \texttt{\tiny{[GM$|$GM]}}
                \texttt{Error: The response should contain the 'guess:' keyword exactly once.} \\
            }
        }
    }
    & & \\ \\

    \theutterance \stepcounter{utterance}  
    & & & \multicolumn{2}{p{0.3\linewidth}}{
        \cellcolor[rgb]{0.9,0.9,0.9}{
            \makecell[{{p{\linewidth}}}]{
                \texttt{\tiny{[GM$|$GM]}}
                \texttt{game\_result = ABORT} \\
            }
        }
    }
    & & \\ \\

    \theutterance \stepcounter{utterance}  
    & & \multicolumn{4}{p{0.6\linewidth}}{
        \cellcolor[rgb]{0.95,0.95,0.95}{
            \makecell[{{p{\linewidth}}}]{
                \texttt{\tiny{[P1$\leftrightarrow$P]}}
                \texttt{Post{-}episode Reflexion (not sent to GM):} \\
\texttt{New plan: The last attempt was aborted because the response contained a long, unstructured paragraph of reasoning instead of the required two{-}line format. I failed to separate the explanation from the guess and did not explicitly use the "guess:" tag on a new line. Next time, I will strictly output exactly two lines: the first line starting with "explanation:" containing the reasoning that combines the clue "tithing" and the feedback (identifying the word as "tithing" based on the clue and letter positions), and the second line starting with "guess:" followed immediately by the word "tithing", ensuring no extra text, newlines, or conversational filler is included.} \\
            }
        }
    }
    & \\ \\

\end{supertabular}
}

\end{document}
